%&preamble/precompile

\title{Incorporating Uncertainty into Reinforcement Learning through Gaussian Processes}
\author{Markus Kaiser}
\subtitle{Master's Thesis}
\author{\href{mailto:markus.kaiser@in.tum.de}{Markus Kaiser}}
\titlegraphic{\parbox{.95\textwidth}{%
    \raisebox{-.5\height}{\includegraphics[height=1.5cm,width=.33\linewidth,keepaspectratio]{tum}}%
    \hfill%
    \raisebox{-.5\height}{\includegraphics[height=1.5cm,width=.33\linewidth,keepaspectratio]{kth}}%
    \hfill%
    \raisebox{-.5\height}{\includegraphics[height=1.5cm,width=.33\linewidth,keepaspectratio]{siemens}}%
}}
\date{4. July 2016}

% \includeonlyframes{solution_psop,solution_reward}

\begin{document}

\begin{frame}[plain, noframenumbering]
    \titlepage
\end{frame}

\begin{frame}[noframenumbering]
    \frametitle{Outline}

    \tableofcontents
\end{frame}

\section{The Bicycle Benchmark}
\begin{frame}[label=bicycle_drive]
    \frametitle{Bicycle Benchmark -- Driving}

    \begin{columns}
        \begin{column}[c]{0.5\textwidth}
            \centering
            \includestandalone{figures/bicycle_from_above}
        \end{column}
        \begin{column}[c]{0.5\textwidth}
            \centering
            \vspace{1cm}
            \uncover<5-> {
                \includestandalonewithpath{figures/bicycle_drive_goal}
            }
        \end{column}
    \end{columns}
\end{frame}

\begin{frame}[label=bicycle_balance]
    \frametitle{Bicycle Benchmark -- Balancing}

    \vfill
    \begin{columns}
        \begin{column}{0.45\textwidth}
            \centering
            \includestandalone{figures/bicycle_from_behind}
        \end{column}
        \begin{column}{0.55\textwidth}
            \centering
            \uncover<4-> {
                \includestandalonewithpath{figures/bicycle_drive_balance}
            }
        \end{column}
    \end{columns}
\end{frame}

\begin{frame}[label=bicycle_summary]
    \frametitle{Bicycle Benchmark -- Summary}

    \centering
    \begin{columns}[t]
        \begin{column}{.4\textwidth}
            \begin{table}
                \caption{State Variables}
                \begin{tabularx}{\textwidth}{cX}
                    \toprule
                    Notation & Description \\
                    \midrule
                    $\theta$ & Steering Angle \\
                    $\dot{\theta}$ &  \\
                    $\omega$ & Leaning Angle \\
                    $\dot{\omega}$ &  \\
                    $x$ & Front tyre position \\
                    $y$ & Front tyre position \\
                    $\psi$ & Bicycle orientation \\
                    \bottomrule
                \end{tabularx}
            \end{table}
        \end{column}
        \begin{column}{.4\textwidth}
            \begin{table}
                \caption{Actions}
                \begin{tabularx}{\textwidth}{cX}
                    \toprule
                    Notation & Description \\
                    \midrule
                    $d$ & Lean distance \\
                    $T$ & Steering force \\
                    \bottomrule
                \end{tabularx}
            \end{table}
        \end{column}
    \end{columns}
\end{frame}

\section{Gaussian Processes}
\begin{frame}[label=gp_non_parametric_models]
    \frametitle{Bayesian Non-Parametric Models}

    \begin{overlayarea}{\textwidth}{8\baselineskip}
        \only<1> {
            \begin{block}{Data}
                \begin{itemize}
                    \item Observe a \structure{noisy} data set $\Set{\left( \mat{x_i}, \mat{y_i} \right)}$
                    \item Assume that $\mat{y_i} = f(\mat{x_i}) + \epsilon$
                \end{itemize}
            \end{block}
        }
        \only<2-3> {
            \begin{block}{Parametric Models}
                \begin{itemize}
                    \item Assume \structure{structure} about $f$, such as $f(\mat{x}) = \mat{W}\mat{x}$
                    \item Find a \structure{single} $\mat{W}$ which is optimal w.r.t.\ some criterion
                \end{itemize}
            \end{block}
        }
        \only<4> {
            \begin{block}{Bayesian Parametric Models}
                \begin{itemize}
                    \item Assume \structure{structure} about $f$, such as $f(\mat{x}) = \mat{W}\mat{x}$
                    \item Find a \structure{distribution} $\Prob{\mat{W} \given \mat{X}, \mat{y}}$ of plausible parameters
                \end{itemize}
            \end{block}
        }
        \only<5-> {
            \begin{block}{Bayesian Non-Parametric Models}
                \begin{itemize}
                    \item $f$ can be an \structure{arbitrary function}
                    \item Find a \structure{distribution} $\Prob{f \given \mat{X}, \mat{y}}$ over functions
                \end{itemize}
            \end{block}
        }
    \end{overlayarea}
    \vfill
    \centering
    \includestandalonewithpath{figures/non_parametric_models}
\end{frame}

\begin{frame}[label=gp_definition]
    \frametitle{Gaussian Processes}

    \begin{definition}[Gaussian Process]
        A \structure{Gaussian Process (GP)} is a collection of random variables $\Set{\rv{F_x}}$, any finite subset of which has a joint Gaussian distribution.
    \end{definition}
    \begin{itemize}
        \item Extension of Gaussians to (infinite) function spaces
        \item $\rv{F_x}$ models the function value $f(x)$
    \end{itemize}
    \vfill
    \pause
    \begin{block}{Mean and Kernel Functions}
        A GP is \structure{completely determined} by two functions.
        \medskip
        \begin{description}[\quad Kernel function]
            \item[Mean function] \makebox[\widthof{$\K(x, x^\prime)$}][r]{$\mu_f(x)$}${} = \Moment{\E}{f(x)}$
            \smallskip
            \item[Kernel function] $\K(x, x^\prime) = \Moment{\cov}{f(x), f(x^\prime)}$
        \end{description}
    \end{block}
    \begin{itemize}
        \item The kernel encodes the \structure{prior assumptions} about the function
    \end{itemize}
\end{frame}

\begin{frame}<1-5,7>[label=gp_hyperparameters]
    \frametitle{Hyperparameters and Posterior}

    \centering
    \includestandalonewithpath{figures/gp_hyperparameters}
\end{frame}

\begin{frame}[label=gp_sparse]
    \frametitle{Sparse Gaussian Processes}

    \begin{overlayarea}{\textwidth}{7\baselineskip}
        \only<-2> {
            \begin{block}{Full Gaussian Processes}
                \begin{itemize}
                    \item Calculation of the posterior in $\Oh(N^3)$
                    \item Not feasible for large data sets
                \end{itemize}
            \end{block}
        }
        \only<3-> {
            \begin{block}{Sparse Gaussian Processes}
                \begin{itemize}
                    \item Find a set of $M$ \structure{Pseudo Inputs}
                    \item Calculation of the posterior in $\Oh(NM^2)$
                \end{itemize}
            \end{block}
        }
    \end{overlayarea}
    \vfill
    \centering
    \includestandalonewithpath{figures/spgp}
\end{frame}

\section{Model-Uncertainties in Reinforcement Learning}
\begin{frame}<1-5>[label=solution_architecture]
    \frametitle{Solution Architecture}

    \centering
    \includestandalone{figures/agent_environment}
\end{frame}

\begin{frame}[label=solution_psop]
    \frametitle{PSO-Policy}

    \begin{definition}[PSO-Policy]
        \label{def:psop}
        The \structure{Particle Swarm Optimization-Policy (PSO-P)} chooses actions via optimization of the expected accumulated reward.
        \begin{align}
            \pi_{\text{PSO-P}}(\mat{s}) &\coloneqq \mat{a^\ast}_0\text{, where} \\
            \mat{a^\ast} &\in \argmax_{\mat{a} \in \Ah^T} \left\{ \Moment*{\E}{\sum_{t=1}^T \gamma^t \Rwd(\mat{s_t}) \given \GP, \mat{s_0} = \mat{s}, \mat{a_0}, \dots, \mat{a_{T-1}}} \right\}
        \end{align}
    \end{definition}
    \begin{itemize}
        \item PSO is a gradient-free heuristic
        \item Directly exploits the transition models
    \end{itemize}
\end{frame}

\againframe<6>{solution_architecture}

\begin{frame}[label=solution_map]
    \frametitle{MAP without Uncertainties}

    \centering
    \includestandalonewithpath{figures/map}
\end{frame}

\begin{frame}[label=solution_map_predictions]
    \frametitle{MAP Predictions}

    \centering
    \includestandalonewithpath{figures/map_predictions}
\end{frame}

\begin{frame}[label=solution_os]
    \frametitle{One-Step Uncertainties}

    \centering
    \includestandalonewithpath{figures/one_step}
\end{frame}

\begin{frame}[label=solution_os_predictions]
    \frametitle{One-Step Predictions}

    \centering
    \includestandalonewithpath{figures/one_step_uncertainties}
\end{frame}

\begin{frame}[label=solution_ms]
    \frametitle{Multi-Step Uncertainties}

    \centering
    \includestandalonewithpath{figures/linearization}
\end{frame}

\begin{frame}[label=solution_ms_predictions]
    \frametitle{Multi-Step Predictions}

    \centering
    \includestandalonewithpath{figures/linearization_predictions}
\end{frame}

\section{Results}
\begin{frame}[label=results_qualitative]
    \frametitle{Qualitative Results}

    \centering
    \includestandalonewithpath{figures/xy_trajectories}
\end{frame}

\begin{frame}[label=results_quantitative]
    \frametitle{Quantitative Results}

    \centering
    \vfill
    \begin{columns}
        \begin{column}[c]{0.5\textwidth}
            \centering
            \uncover<3> {
                \includestandalonewithpath{figures/goal_percentage}
            }
        \end{column}
        \begin{column}[c]{0.5\textwidth}
            \centering
            \begin{tabularx}{\textwidth}{X>{\centering\hsize=.4\hsize}X>{\centering\hsize=.4\hsize}X>{\centering\hsize=.4\hsize\arraybackslash}X}
                \toprule
                Metric & MAP & OS & MS \\
                \midrule
                Trajectories & 9660 & 9660 & 9660
                \uncover<2-> {
                    \\
                    \addlinespace
                    Success Rate & \SI[mode=text]{53.4}{\percent} & \SI[mode=text]{63.8}{\percent} & \textbf{\SI[mode=text,detect-weight]{75.8}{\percent}} \\
                    \addlinespace
                    Time to Goal \\
                    Mean & \textbf{59.9} & 62.0 & 66.5 \\
                    Median & \textbf{60} & \textbf{60} & 63
                }
                \\
                \bottomrule
            \end{tabularx}
            \vspace{2ex}
            \begin{description}[MAP]
                \item[MAP] Deterministic predictions
                \item[OS] One-Step Uncertainties
                \item[MS] Multi-Step Uncertainties
            \end{description}
        \end{column}
    \end{columns}
\end{frame}

\appendix
\section{Additional Material}
\begin{frame}
    \vfill
    \centering
    \begin{beamercolorbox}[sep=8pt,center]{title}
        \usebeamerfont{title}Additional Material\par%
    \end{beamercolorbox}
    \vfill
\end{frame}

\againframe<5-7>{gp_hyperparameters}

\begin{frame}[label=solution_reward]
    \frametitle{Reward Function}

    \begin{definition}[Bicycle Reward Function]
        The \structure{bicycle reward function} is defined as
        \begin{align}
            \RwdBicycle(\mat{s}) &\coloneqq \begin{cases}
                    2 & \text{if goal reached} \\
                    0 & \text{if fallen down} \\
                    c\cdot\Gaussian{\Delta^\psi_{\mat{s}} \given 0, \sigma_{\text{angle}}^2} & \text{otherwise}
                \end{cases}
        \end{align}
    where $\Delta^\psi_{\mat{s}}$ denotes the \structure{angle towards the goal}.
    \end{definition}
    \vspace{2ex}
    \centering
    \includestandalone{figures/saturating_reward}
\end{frame}

\end{document}
